% =====================================================================
% Rozdział 5: Zagadnienia implementacyjne
% =====================================================================

\chapter{Zagadnienia implementacyjne}
\label{chap:implementacja}

\section{Przetwarzanie reprezentacji}
Implementacja LEM MUSI umożliwiać odwzorowanie danych zapisanych w reprezentacji LEM na komponenty modelu semantycznego.
Proces przetwarzania NIE MOŻE zmieniać interpretacji komponentów określonej w Rozdziale 2. Dotyczy to również opisów częściowych (zob. 2.4.4) — brak komponentu MUSI być odwzorowany jako brak informacji o mniejszej precyzji, a nie interpretowany jako błąd. Zasady postępowania z elementami nierozpoznanymi przez implementację są określone w sekcji 7.7.

\section{Reprezentacja wewnętrzna implementacji}
Implementacja MOŻE stosować wewnętrzne struktury danych różne od reprezentacji wymiany. Wewnętrzna reprezentacja NIE MOŻE wpływać na semantykę informacji LEM.

\section{Walidacja danych}
Implementacja MOŻE wykonywać walidację w następujących zakresach:

\begin{itemize}
    \item poprawność składniowa — zgodność zapisu ze strukturą wymaganą przez reprezentację,
    \item zgodność z regułami reprezentacji — poprawność odwzorowania elementów reprezentacji na komponenty modelu (obejmuje ocenę kompletności danych wymaganej przez daną reprezentację),
    \item zgodność z profilem zastosowania — spełnienie wymagań określonego profilu (zob. 4.4.3).
\end{itemize}

Walidacja NIE MOŻE definiować nowej interpretacji informacji.

\section{Obsługa błędów}
Implementacja MUSI rozróżniać kategorie błędów odpowiadające zakresom walidacji z sekcji 5.3:

\begin{itemize}
    \item błędy składni reprezentacji,
    \item błędy zgodności z regułami reprezentacji,
    \item błędy zgodności z profilem.
\end{itemize}

Sposób techniczny sygnalizowania tych kategorii błędów (kody błędów, mechanizmy raportowania itp.) pozostaje poza zakresem modelu LEM.

Dane niepoprawne lub niejednoznaczne NIE MOGĄ być interpretowane jako poprawna informacja LEM — jest to bezpośrednia konsekwencja zasady jednoznaczności interpretacji określonej w sekcji 2.5.1.

\section{Wersjonowanie i kompatybilność implementacji}
Zmiany implementacyjne NIE MOGĄ powodować zmiany znaczenia komponentów modelu LEM.

Implementacje POWINNY umożliwiać identyfikację wersji używanych reprezentacji oraz zachowanie możliwości przetwarzania danych zgodnych z wcześniejszymi wersjami modelu (zob. Rozdział 7).